% !PRINTED 612
\weberpage{612}

If, then, in a group $L$ there occurs a substitution $\lambda$ in which $a = 1$ and $b$ is different from zero, then by the above $L$ contains the whole \gesperrt{cyclic group} $(z, z + b)$ $(b = 0, 1, \ldots n - 1)$.

If $a$ is not equal to 1, we set
\[
1 + a + \cdots + a^{h-1} = \frac{a^h - 1}{a - 1},
\]
and conclude from this that, if $e$ is the least positive exponent for which $a^e \equiv 1\ (\mathrm{mod}\ n)$, then $\lambda$ is of degree $e$.

The period of the substitution $\lambda$ is, unless $a = 1$, an intransitive group. There is a figure $z$ which is not altered by $\lambda$ and its repetitions, which is determined from the congruence $z \equiv az + b\ (\mathrm{mod}\ n)$, and can be denoted by
\[
z_0 \equiv \frac{b}{a - b}\ (\mathrm{mod}\ n)
\]
% [sic] source denominator is "a-b"; fixed-point algebra would give "1-a"
if by $\dfrac{\mu}{\nu}\ (\mathrm{mod}\ n)$ is understood an integer which, by multiplication with the factor $\nu$, becomes congruent to $\mu$ with respect to the modulus $n$.

The cyclic group $(z, z + b)$ is transitive and of degree $n$. It can be shown that every transitive linear group $(\mathrm{mod}\ n)$ must contain the cyclic group.

Let $g$ be a primitive root of $n$ and $\alpha$ the index of $a$, i.\,e. $g^\alpha \equiv a\ (\mathrm{mod}\ n)$ (\S.\,136). For brevity we call $\alpha$, for the moment, also the \gesperrt{index of the substitution} $\lambda = (z, az + b)$. The index may be assumed to be less than $n - 1$ and equal to or greater than zero. If it is equal to zero, then $\lambda$ is either the identical substitution, or it belongs to the cyclic group. By (4) the rule holds that the index of a compounded substitution is equal to the sum of the indices of the components.

From this it follows that all the indices of the substitutions of a linear group $L$ are multiples of the least positive one among them, and that consequently, if $\alpha_0$ is this least positive index and $a^{\alpha_0} = a_0$ is set, all the substitutions of $L$ can be represented in the form
% [sic] The source prints a^{\alpha_0}=a_0 here; the preceding definition uses g^{\alpha}=a. Retained source-faithfully.
\[
\lambda = (z, a_0^h z + b)
\]
in which $a_0$ is held fixed, and only $h$ and $b$ run through certain sequences of numbers.

A substitution of the form $\lambda_0 = (z, a_0 z + b_0)$ must occur in the group $L$.
